\documentclass[acmsmall,screen,review,anonymous]{acmart}

\setcopyright{none}
\settopmatter{printacmref=true}
\renewcommand\footnotetextcopyrightpermission[1]{}

\acmConference[miniKanren 2026]{The 2026 miniKanren and Relational Programming Workshop}{August 24, 2026}{Indianapolis, Indiana, USA}
\acmYear{2026}
\copyrightyear{2026}

\title{Proflog: A Tableau-Based Logic Programming Kernel in miniKanren}

\begin{document}

\begin{abstract}
We report on an implementation of Melvin Fitting's Proflog, a logic programming
language based on the tableau proof procedure.  The system extends a first-order
tableau prover with equality, written in Clojure's \texttt{core.logic}, with
Fitting's Procedure Call rule.  Formulae thereby acquire a programmatic
character: evaluating a program is a matter of closing the relevant tableau.
Following alphaLeanTAP, successful evaluation also produces an explicit proof
object.  We describe the proof kernel, the surface-to-kernel translation, and
worked executions of Fitting's reference programs.
\end{abstract}

\ccsdesc[500]{Theory of computation~Logic and verification}
\ccsdesc[300]{Theory of computation~Automated reasoning}
\ccsdesc[300]{Software and its engineering~Constraint and logic languages}

\keywords{miniKanren, logic programming, tableaux, Proflog, proof search}

\maketitle

\section{Introduction}

Fitting proposed Proflog as an exploration in alternative bases for logic
programming~\cite{fitting1994tableaux}.  Prolog operationalizes Horn clauses by
resolution.  Proflog instead manipulates formulae of first-order logic with
equality by a tableau proof procedure, augmented with a mechanism for recursive
procedure calls.  The core tableau procedure is standard: branch closure,
alpha/beta decomposition of connectives, gamma/delta treatment of quantifiers,
and equality over a weak Herbrand setting.  The programming step is the
Procedure Call rule.  During an attempt to close a tableau, a saved literal
whose predicate is defined by the program may open a subsidiary tableau for the
corresponding clause body.

This changes the usual programming/proving boundary.  In Proflog, a program is
not a collection of Horn clauses with procedural negation as failure.  It is a
finite collection of first-order definitions, at most one per relation symbol,
whose queries are interpreted by two proof searches: a query succeeds when the
negated query closes, fails when the query itself closes, and is unresolved when
neither direction closes within the admitted search.  This mirrors Fitting's
supervaluation semantics and makes some operational warnings visible, especially
the warning that factoring a formula into an auxiliary relation may change the
program.

The implementation reported here is written in Clojure and \texttt{core.logic}
in the miniKanren tradition~\cite{byrd2009relational,corelogic}.  Its immediate
prover lineage is alphaLeanTAP, a relational tableau prover for first-order
classical logic~\cite{near2008alphaleantap}, and a Clojure/core.logic.nominal
implementation reference by Amin~\cite{aminLeantap}.  Proflog extends that
lineage in the direction Fitting sketched: first-order proof search becomes a
program evaluator, and program evaluation returns proof objects whose step tags
record the tableau work performed.

\section{Implementation}

\subsection{Proof Kernel}

The kernel is a relation over proof states.  Its central entry point is
\texttt{prove-stateo}; public calls use \texttt{kernel/prove} for direct
formulae and \texttt{kernel/prove-program} for program-bearing formulae.  A
branch state contains the focused formula, pending formulae, saved literals,
the lexical environment for bound variables, gamma-introduced proof variables,
delta parameters, an explicit equality substitution, delayed disequalities, the
compiled program, fuel, and the proof term being constructed.

The equality design is deliberately explicit.  Positive equality extends a
free-constructor substitution; negative equality is retained in a symbolic
disequality store until later bindings force closure.  This keeps equality,
disequality, and procedure-call interaction visible in proof objects.  A small
closure proof, for example, has the shape:

\begin{verbatim}
(conj (savefml (close)))
\end{verbatim}

An equality/disequality proof exposes the order of work:

\begin{verbatim}
(conj (neq-store (eq-step (eq-bind) (neq-close))))
\end{verbatim}

Procedure calls add \texttt{pos-call} and \texttt{neg-call} proof steps.  A
positive call closes the relation body.  A negative call closes the normalized
negation of that body.  Both calls run under the same kernel relation rather
than through host-language evaluation of the source program.

\subsection{Surface Syntax and Descent}

The implemented frontend is a Clojure-readable prefix DSL that preserves
Fitting's clause organization.  The intended handwritten notation is:

\begin{verbatim}
p(x) :- x = a.

only-zero(x) := forall y. (x != y or y = zero).
zero-only(x) :- only-zero(x).
\end{verbatim}

The frontend form uses \texttt{|-} for a real Proflog clause and \texttt{:=}
for a frontend-only abbreviation:

\begin{verbatim}
(def peano-language
  (pf/language
    (constants zero)
    (functions (s 1))
    (relations (zero-only 1))))

(def zero-only-program
  (pf/proflog peano-language
    (:= (only-zero x)
      (forall [y]
        (or (!= x y)
            (= y zero))))
    (|- (zero-only x)
        (only-zero x))))
\end{verbatim}

After macro expansion and compilation, the kernel sees tagged formula data.  The
simple query \texttt{p(a)} descends schematically to:

\begin{verbatim}
(pos (app p (app a)))
\end{verbatim}

The inlined \texttt{zero-only/1} body descends schematically to:

\begin{verbatim}
(forall
  (tie y
    (or
      (neq (var x) (var y))
      (eq (var y) (app zero)))))
\end{verbatim}

The public query API then probes the kernel in two directions:

\begin{verbatim}
(query/query-succeeds program query n fuel)
(query/query-fails    program query n fuel)
\end{verbatim}

\texttt{query-succeeds} asks the kernel to close the negated query;
\texttt{query-fails} asks it to close the query itself.  Open-answer examples
use \texttt{pf/run}, which binds visible answer variables in the style of
miniKanren's \texttt{run}, translates the query to the same backend formula,
and exports bindings from the proof-backed answer layer.

\subsection{Fitting's Reference Programs}

Fitting's first reference program defines even and odd natural numbers over the
language with constant \texttt{zero} and unary successor \texttt{s}:

\begin{verbatim}
even(x) :- x = zero
        or exists y. (x = s(y) and odd(y)).

odd(x) :- forall y. (even(y) -> x != y).
\end{verbatim}

The implemented frontend form is:

\begin{verbatim}
(pf/proflog p1-language
  (|- (even x)
    (or (= x zero)
        (exists [y]
          (and (= x (s y))
               (odd y)))))
  (|- (odd x)
    (forall [y]
      (implies (even y)
               (!= x y)))))
\end{verbatim}

The second reference program is a one-or-two-token Nim game.  A position wins
when there is a legal move to a position that does not win:

\begin{verbatim}
win(x) :- exists y.
            ((x = s(y) or x = s(s(y))) and not win(y)).
\end{verbatim}

The implementation keeps the move test inlined, as in Fitting's presentation:

\begin{verbatim}
(pf/proflog p2-language
  (|- (win x)
    (exists [y]
      (and (or (= x (s y))
               (= x (s (s y))))
           (not (win y))))))
\end{verbatim}

Table~\ref{tab:fitting-results} summarizes the focused Fitting suite results.
The suite was run as a kernel-level test gate and completed with six tests,
eighty-one assertions, and no failures or errors.  It covers forward success,
forward refutation, answer synthesis, partial synthesis, and unresolved
classification.

\begin{table}[h]
  \caption{Representative Fitting-program outcomes}
  \label{tab:fitting-results}
  \begin{tabular}{lll}
    \toprule
    Case & Mode & Outcome \\
    \midrule
    \texttt{even(0)} & forward success & succeeds \\
    \texttt{odd(s(0))} & forward success & succeeds \\
    \texttt{odd(0)} & forward refutation & fails \\
    \texttt{win(4)} & forward success & succeeds \\
    \texttt{win(3)} & forward refutation & fails \\
    \texttt{move(1,0)} & auxiliary forward success & succeeds \\
    \texttt{move(0,1)} & auxiliary forward refutation & fails \\
    factored \texttt{win(1)} & status classification & unresolved \\
    \texttt{append(x,y,[a,b,c])} & answer synthesis & target found \\
    \texttt{reverse([a,b,c,a],cons(a,r))} & partial synthesis & target found \\
    \bottomrule
  \end{tabular}
\end{table}

The factored Nim case is included because it demonstrates that the evaluator is
not merely computing the Nim recurrence in the host language.  If the program
introduces an auxiliary relation

\begin{verbatim}
win(x) :- exists y. (move(x, y) and not win(y)).
move(x, y) :- x = s(y) or x = s(s(y)).
\end{verbatim}

then \texttt{move/2} itself is decidable on ground calls, but \texttt{win(1)}
remains unresolved in the tested kernel setting.  This is the operational
boundary Fitting warned about: a logically tempting refactoring changes how
procedure calls are admitted and discharged.

\subsection{Test Strategy}

The implementation separates routine regression from slow semantic probes.  The
Fitting suite disables older host-side overlays while running, so its promoted
rows cannot pass by fixture-specific evaluation.  Representative proof objects
are inspected for ordinary kernel and answer-overlay evidence, including
procedure-call tags and the absence of diagnostic sidecar settlement tags.  The
same suite also includes finite-domain examples, list programs, and group
verifier formulae to check that the proof kernel and proof profiles are not
tailored only to the two paper examples.

\section{Discussion}

Proflog is useful to the miniKanren setting because it gives another concrete
point in the space between logic and computation.  The program text is not a
relational interpreter written in a host language, nor is it a Prolog program
embedded in Clojure.  It is a collection of first-order formulae whose
operational content comes from tableau closure plus a procedure-call rule.  That
makes proof search, program evaluation, and proof-object construction part of
one relational computation.

The implementation is also a cautionary example.  Fitting's semantics are not
Prolog's semantics with a different syntax.  Query failure is proved, not
assumed by search exhaustion.  Undefined and unresolved cases are first-class
outcomes.  Refactorings that would be harmless in a purely extensional reading
can change operational behavior when they introduce auxiliary procedure calls.
These facts show up in tests rather than only in prose.

The current system is not presented as an efficient logic programming language.
Some deeper recursive examples remain search-control frontiers, and several
successful answer and partial-synthesis tests are intentionally kept out of the
default fast path.  The contribution is instead the executable correspondence:
Fitting's theoretical Proflog can be mechanized as a proof-producing
miniKanren-style kernel, and its non-trivial examples can be evaluated through
that kernel with correctness results visible at the proof boundary.

\begin{acks}
Acknowledgments are omitted for anonymous review.
\end{acks}


\bibliographystyle{ACM-Reference-Format}
\bibliography{references}

\end{document}
